% Figure: relative-motion / river-crossing vector triangle. The boat's
% velocity relative to the water adds to the water's velocity relative
% to the ground to give the boat's velocity relative to the ground.
% Two steering choices are shown: heading straight across (swept
% downstream) and heading upstream to land directly opposite.
\begin{figure}[htbp]
\centering
\begin{tikzpicture}[
  vb/.style={draw=engblue, -{Stealth}, very thick},
  vw/.style={draw=engamber, -{Stealth}, very thick},
  vg/.style={draw=englime!60!black, -{Stealth}, very thick},
  lbl/.style={font=\scriptsize},
]
% banks
\draw[draw=enggray, thick] (-0.4,0) -- (8.5,0);
\draw[draw=enggray, thick] (-0.4,4) -- (8.5,4);
\fill[engblue!6] (-0.4,0) rectangle (8.5,4);
\node[lbl, enggray, anchor=west] at (-0.4,-0.3) {near bank};
\node[lbl, enggray, anchor=west] at (-0.4,4.3) {far bank};
% current direction arrow (downstream, to the right)
\node[lbl, engamber, anchor=west] at (5.6,2) {current $\vec{v}_{w}$};
\draw[vw] (5.4,2) -- ++(1.6,0);

% Case A: head straight across, swept downstream.
\coordinate (A0) at (0.6,0);
% boat rel water: straight up, magnitude 3 (scaled: 3 -> 3 units)
\draw[vb] (A0) -- ++(0,3) node[midway, left, lbl, engblue] {$\vec{v}_{b/w}$};
% water rel ground: to the right, magnitude 1.5
\draw[vw] (A0) ++(0,3) -- ++(1.5,0);
% boat rel ground: resultant from A0 to the tip
\draw[vg] (A0) -- ++(1.5,3) node[midway, right, lbl, englime!50!black] {$\vec{v}_{b/g}$};
\node[lbl, anchor=north] at (0.6,-0.1) {(a) point straight across};

% Case B: head upstream so resultant is straight across.
\coordinate (B0) at (4.3,0);
% boat rel water aimed up-left so that across-component lands opposite:
% magnitude 3, with x-component -1.5 to cancel current; y = sqrt(3^2-1.5^2)=2.598
\draw[vb] (B0) -- ++(-1.5,2.598) node[midway, left, lbl, engblue] {$\vec{v}_{b/w}$};
% water rel ground to the right, magnitude 1.5
\draw[vw] (B0) ++(-1.5,2.598) -- ++(1.5,0);
% boat rel ground: straight up, magnitude 2.598
\draw[vg] (B0) -- ++(0,2.598) node[midway, right, lbl, englime!50!black] {$\vec{v}_{b/g}$};
\node[lbl, anchor=north] at (4.3,-0.1) {(b) angle upstream};
\end{tikzpicture}
\caption[Relative-motion vector triangle for a river crossing]{The
two steering choices for crossing a current, each a closed vector
triangle of the relative-motion law $\vec{v}_{b/g}=\vec{v}_{b/w}+
\vec{v}_{w/g}$. In (a) the boat points straight across, so its velocity
relative to the water (blue) is perpendicular to the banks; the current
(amber) adds a downstream component and the velocity relative to the
ground (green) carries the boat across on a slanted track, landing
downstream of the start. The crossing time is shortest here because the
full speed counts toward the width. In (b) the boat angles upstream by
just enough that the current cancels the sideways drift, so the
ground velocity points straight across and the landing is directly
opposite; the price is a smaller across-component and a longer crossing
time. The triangle is the same law that governs an aircraft crabbing
into a crosswind.}
\label{fig:vol03ch04:river-crossing}
\end{figure}
